\Askhsh{37009}{2}
\begin{ylh}{1}
    \item 3.1. Είδη και στοιχεία τριγώνων
    \item 4.2. Τέμνουσα δύο ευθειών - Ευκλείδειο αίτημα
    \item 4.6. Άθροισμα γωνιών τριγώνου.
\end{ylh}
% ===================================================================
%  Αρχείο: 37009.tex (Fragment)
%  Αυτόματη μετατροπή μέσω doc2latex.py
% ===================================================================

ΘΕΜΑ 2

Δίνεται ορθογώνιο τρίγωνο ΑΒΓ ($\widehat{Α} = 90^{0}$) και ΑΔ η διχοτόμος της γωνίας $\widehat{Α}$. Από το σημείο Δ φέρουμε την παράλληλη προς την ΑΒ που τέμνει την ΑΓ στο Ε.

\begin{alist}
\item Να αποδείξετε ότι το τρίγωνο ΕΔΓ είναι ορθογώνιο. \monades{9}

\item Να υπολογίσετε τη γωνία Α$\widehat{\Delta}$Ε. \monades{9}

\item Αν η γωνία $\widehat{Β}$ είναι 20\textsuperscript{ο} μεγαλύτερη της γωνίας $\widehat{\Gamma}$, να υπολογίσετε τη γωνία $Ε\widehat{\Delta}\Gamma$.

\monades{7}
\end{alist}
